% EDN entry template. Transfer an entry from reports/edn.md: copy this file to
% edn/NN-short-slug.tex, fill it in, and add
% \input{edn/NN-short-slug} to the entry list in edn.tex.
% Brace every value. Delete optional lines you do not need.
\ednentry{
  id = {EDN-NN},
  title = {Short name of the decision},
  date = {YYYY-MM-DD},                   % optional
  % M1: Project Inception | M2: Reproducibility | M3: Quality Assurance |
  % M4: Model Deployment | M5: Model Packaging | M6: Monitoring | Other
  milestone = {M1: Project Inception},
  topic = {e.g. Testing Strategy},       % optional
  participants = {All team members},     % or the names of who decided
  decision = {What we decided.},
  rationale = {Why; alternatives with their pros and cons, constraints, trade-offs, evidence.},
  % Select all that apply: No AI involvement; Information seeking;
  % Alternative generation; Alternative assessment; Recommendation;
  % Solution generation; Other (describe).
  ai = {No AI involvement.},
  % Only if AI contributed. Select one: Accepted; Accepted with modifications;
  % Rejected; Used as input for further analysis; Other (describe).
  response = {},
  assessment = {},                       % why we accepted/modified/rejected it
  % Optional: link to the conversation (say which part), or quote it:
  % \ednprompt{...} \ednresponse{...}
  % Review shared content first; never share personal or sensitive data.
  aievidence = {},
  evidence = {},                         % optional: issues, PRs, commits, runs
}
